% Vietnamese translation for OI Public Library
% Source-File: IOI中国国家候选队论文/2003/雷环中/雷环中.ppt
\documentclass[11pt]{article}
\usepackage{oipl-vn}

\title{Bản trình chiếu: Bài toán dạng nộp kết quả}
\author{Lei Huanzhong}
\date{2003}

\begin{document}
\maketitle

\origtitle{On output-only problems: presentation notes}
\sourcepath{IOI China candidate team papers / 2003 / Lei Huanzhong / Lei Huanzhong.ppt}

\section{Ý chính của bài trình bày}

Bộ slide là bản trình bày ngắn đi kèm bài viết về bài toán dạng nộp kết quả.
Nó không cố trình bày một thuật toán tổng quát, mà dùng hai ví dụ cụ thể để
nhấn mạnh đặc trưng của kiểu bài này: dữ liệu đầu vào đã biết trước, kết quả
cuối cùng quan trọng hơn quá trình sinh kết quả, và vì vậy người giải có thể
kết hợp phân tích thủ công, chương trình phụ trợ ngắn, thống kê và thử nghiệm.

Các phần chữ khôi phục được từ PowerPoint cho thấy trọng tâm của slide là
\textit{Crack the Code} của Baltic OI 2001 và \textit{Tetris} của NOI 2002.
Những nội dung rộng hơn như \textit{Birthday Party}, \textit{Xor} hay
\textit{Double Crypt} đã được trình bày đầy đủ hơn trong bài viết cùng tác
giả; bản ghi chú này chỉ nối lại các luận điểm có mặt trong deck.

\section{\textit{Crack the Code}}

\subsection{Dữ liệu của bài}

Slide mở đầu phần này nhắc tới ba tệp:
\[
  \texttt{cra.out},\qquad \texttt{cra.txt},\qquad \texttt{cra.in}.
\]
Tệp \texttt{cra.txt} là văn bản gốc tiếng Anh; \texttt{cra.out} là văn bản
cần khôi phục; \texttt{cra.in} là kết quả thu được sau khi áp dụng phép mã hóa
trên \texttt{cra.out}. Từ \texttt{cra.txt} và \texttt{cra.in}, người giải phải
tìm ra \texttt{cra.out}.

Ký hiệu chính trong slide là phép
\[
  \operatorname{succ}(C,0)=C,\qquad
  \operatorname{succ}(C,n)=\operatorname{succ}(\operatorname{succ}(C,n-1),1).
\]
Nói cách khác, \(\operatorname{succ}(C,n)\) là ký tự nhận được sau khi dịch
ký tự \(C\) tiến \(n\) bước trong bảng chữ cái. Khi xử lý ký tự thứ \(i\), hệ
mã dùng một trong mười tham số \(a_1,a_2,\ldots,a_{10}\), lặp lại theo chu kỳ:
\[
  c_i=\operatorname{succ}\bigl(c_i,\ a_{(i-1)\bmod 10}+1\bigr).
\]
Vì chỉ có mười tham số, một khi đoán đúng vài cặp chữ tương ứng giữa bản mã và
bản rõ, phần còn lại có thể được lấp dần.

\subsection{Phá mã bằng nhận dạng văn bản}

Deck chép lại dữ liệu thứ nhất. Văn bản rõ bắt đầu bằng đoạn nói về Alfred
Tarski:
\begin{quote}
ALFRED TARSKI WAS ONE OF THE GREATEST MATHEMATICIANS, LOGICIANS AND
PHILOSOPHERS OF THE PASSING CENTURY, AND ONE OF THE GREATEST LOGICIANS OF
ALL TIME.
\end{quote}
Trong bản mã tương ứng xuất hiện cụm \texttt{EQMDY 1943}. Dựa vào ngữ cảnh
tiểu sử, cụm này gần như chắc chắn là \texttt{AFTER 1943}. Chỉ riêng giả thiết
này đã cho:
\[
  a_6=4,\qquad a_7=11,\qquad a_8=19,\qquad a_9=25,\qquad a_{10}=7.
\]

Sau khi thay các tham số đã biết, nhiều cụm khác hiện ra dưới dạng gần đọc
được:
\[
\begin{array}{rcl}
\texttt{THG NXHPARD UNKBBHMITY} &\to& \texttt{THE HARVARD UNIVERSITY},\\
\texttt{UNIVETYFJS OF CALKLLHHIA} &\to& \texttt{UNIVERSITY OF CALIFORNIA}.
\end{array}
\]
Những cụm dài như vậy cho thêm các vị trí chưa biết của chu kỳ \(10\), và quá
trình khôi phục trở thành một bài điền tham số tương đối ngắn. Slide dùng ví
dụ này để chỉ ra vai trò của trực giác con người: một từ ngắn hoặc một cụm
tên riêng có thể đem lại nhiều thông tin hơn một chương trình tổng quát chạy
mù.

\subsection{Thống kê tần suất}

Slide cũng gợi ý cách dùng phân bố tần suất. Với mỗi lớp vị trí
\[
  i,\ i+10,\ i+20,\ldots,
\]
ta lấy các chữ trong \texttt{cra.in} thuộc lớp đó, thử mọi dịch chuyển
\(0,\ldots,25\), rồi so sánh phân bố chữ cái thu được với phân bố chữ cái của
văn bản tiếng Anh hoặc của \texttt{cra.txt}. Khi dữ liệu đủ dài, khoảng cách
giữa hai phân bố giúp chọn tham số đúng.

Điểm đáng chú ý là slide không xem thống kê là lời giải thay thế hoàn toàn
cho phân tích thủ công. Khi văn bản ngắn, tần suất dao động mạnh; khi văn bản
có tên riêng hoặc thuật ngữ chuyên ngành, phân bố cũng lệch. Cách làm thực tế
là dùng thống kê để sinh vài ứng viên tốt, rồi kiểm tra bằng mắt và bằng các
từ khóa dễ nhận ra.

\section{\textit{Tetris}}

Phần thứ hai của deck chuyển sang \textit{Tetris} của NOI 2002. Các slide
liệt kê nhiều tệp dữ liệu với kích thước khác nhau:
\[
\begin{array}{c|c}
\text{tệp} & N\\ \hline
\texttt{Tetris4.in} & 16\\
\texttt{Tetris5.in} & 47\\
\texttt{Tetris6.in} & 97\\
\texttt{Tetris7.in} & 100\\
\texttt{Tetris8.in} & 246\\
\texttt{Tetris9.in} & 574\\
\texttt{Tetris10.in} & 1202
\end{array}
\]
Một dòng dữ liệu mẫu trong slide có dạng
\[
  0,2,4,6,6,8,10,12,12,14,16,18,18,20,22,24,\ldots
\]
cho thấy các cột có cấu trúc khá đều. Đây là thông tin quan trọng đối với bài
nộp kết quả: ta không cần tìm một chiến lược duy nhất tối ưu cho mọi bàn chơi
tưởng tượng, mà cần tận dụng dạng của các bàn chơi thật trong bộ dữ liệu.

Theo bài viết đi kèm, một cách làm nhanh là viết chương trình sinh thao tác
theo kinh nghiệm, chạy nhiều lần rồi chọn kết quả tốt. Với dữ liệu nhỏ hoặc
có quy luật rõ, có thể tính tay hoặc dùng quy tắc xây dựng cố định; với dữ
liệu lớn, ngẫu nhiên hóa và hàm đánh giá độ phẳng của mặt bàn thường hiệu
quả hơn. Các con số như \(50\), \(97\), \(200\) xuất hiện trong phần chữ trích
được của slide nhiều khả năng là tham số thử nghiệm hoặc số bước trung gian,
vì vậy bản ghi chú này không biến chúng thành phát biểu thuật toán cụ thể.

\section{Bài học về chiến lược}

Hai ví dụ trong deck bổ sung cho nhau. \textit{Crack the Code} nhấn mạnh phân
tích dữ liệu: đọc dữ liệu thật, đoán cấu trúc, dùng tri thức ngôn ngữ, rồi để
chương trình kiểm tra. \textit{Tetris} nhấn mạnh thử nghiệm: viết bộ sinh kết
quả đủ nhanh, điều chỉnh các quy tắc đánh giá, chạy nhiều lần và giữ lời giải
tốt.

Điểm chung là cách đánh giá hiệu quả đã thay đổi. Trong bài nộp chương trình,
ta thường muốn một thuật toán đúng cho mọi dữ liệu hợp lệ. Trong bài nộp kết
quả, ta muốn tệp trả lời đúng cho các dữ liệu đã cho, trong thời gian thi hữu
hạn. Vì thế một lời giải có thể gồm nhiều mảnh: vài dòng phân tích thủ công,
một chương trình thống kê, một chương trình kiểm chứng, và một số quyết định
chiến lược về việc dữ liệu nào nên làm trước.

\end{document}
